\begin{question}
    Determine os autovalores, os autovetores e os autoespaços da matriz
    \[
        B=
        \begin{pmatrix}
            4 & 1 & 1\\
            1 & 4 & 1\\
            1 & 1 & 4
        \end{pmatrix}
    \]
\end{question}

\begin{resolution}{Exemplo 4: autovalores}
    O polinômio característico é
    \begin{align*}
        p(\lambda)
        &=\det(B-\lambda I)\\
        &=
        \begin{vmatrix}
            4-\lambda & 1 & 1\\
            1 & 4-\lambda & 1\\
            1 & 1 & 4-\lambda
        \end{vmatrix}\\
        &=(4-\lambda)^3-3(4-\lambda)+2\\
        &=(6-\lambda)(3-\lambda)^2
    \end{align*}

    Portanto
    \[
        \lambda_{1}=6
        \qquad
        \lambda_{2}=3
    \]
\end{resolution}

\begin{resolution}{Exemplo 4: autovetores para $\lambda=6$}
    Resolvemos o sistema
    \begin{align*}
        (B-6I)\vec{x}&=\vec{0}\\
        \begin{pmatrix}
            -2 & 1 & 1\\
            1 & -2 & 1\\
            1 & 1 & -2
        \end{pmatrix}
        \begin{pmatrix}
            x\\y\\z
        \end{pmatrix}
        &=
        \begin{pmatrix}
            0\\0\\0
        \end{pmatrix}
    \end{align*}

    O sistema equivalente é
    \begin{align*}
        -2x+y+z&=0\\
        x-2y+z&=0\\
        x+y-2z&=0
    \end{align*}

    Subtraindo as equações, obtemos
    \begin{align*}
        x&=y\\
        y&=z
    \end{align*}
\end{resolution}

\begin{resolution}{Exemplo 4: autoespaço para $\lambda=6$}
    Portanto
    \[
        E_{6}=
        \left\{
            \vetortri{t}{t}{t}
            \mid t\in\mathbb{R}
        \right\}
    \]

    Ou seja
    \[
        E_{6}=
        \operatorname{span}
        \left\{\vetortri{1}{1}{1}\right\}
    \]

    Os autovetores são os múltiplos não nulos de $\vetortri{1}{1}{1}$
\end{resolution}

\begin{resolution}{Exemplo 4: autovetores para $\lambda=3$}
    Resolvemos o sistema
    \begin{align*}
        (B-3I)\vec{x}&=\vec{0}\\
        \begin{pmatrix}
            1 & 1 & 1\\
            1 & 1 & 1\\
            1 & 1 & 1
        \end{pmatrix}
        \begin{pmatrix}
            x\\y\\z
        \end{pmatrix}
        &=
        \begin{pmatrix}
            0\\0\\0
        \end{pmatrix}
    \end{align*}

    O sistema equivalente é
    \[
        x+y+z=0
    \]

    Escolhendo $x=s$ e $y=t$, temos
    \[
        z=-s-t
    \]
\end{resolution}

\begin{resolution}{Exemplo 4: autoespaço para $\lambda=3$}
    Portanto
    \[
        E_{3}=
        \left\{
            \vetortri{s}{t}{-s-t}
            \mid s,t\in\mathbb{R}
        \right\}
    \]

    Escrevendo como combinação linear
    \begin{align*}
        E_{3}
        &=
        \left\{
            s\vetortri{1}{0}{-1}
            +
            t\vetortri{0}{1}{-1}
            \mid s,t\in\mathbb{R}
        \right\}\\
        &=
        \operatorname{span}
        \left\{
            \vetortri{1}{0}{-1}
            ,
            \vetortri{0}{1}{-1}
        \right\}
    \end{align*}

    Logo
    \[
        \dim(E_{3})=2
    \]
\end{resolution}

\begin{resolution}{Exemplo 4: resposta}
    Os autovalores são
    \[
        \lambda_{1}=6
        \qquad
        \lambda_{2}=3
    \]

    Os autoespaços são
    \begin{align*}
        E_{6}&=
        \operatorname{span}
        \left\{\vetortri{1}{1}{1}\right\}\\
        E_{3}&=
        \operatorname{span}
        \left\{
            \vetortri{1}{0}{-1}
            ,
            \vetortri{0}{1}{-1}
        \right\}
    \end{align*}

    O autoespaço $E_{3}$ possui dimensão dois
\end{resolution}